\documentclass[]{book}
\usepackage{lmodern}
\usepackage{amssymb,amsmath}
\usepackage{ifxetex,ifluatex}
\usepackage{fixltx2e} % provides \textsubscript
\ifnum 0\ifxetex 1\fi\ifluatex 1\fi=0 % if pdftex
  \usepackage[T1]{fontenc}
  \usepackage[utf8]{inputenc}
\else % if luatex or xelatex
  \ifxetex
    \usepackage{mathspec}
  \else
    \usepackage{fontspec}
  \fi
  \defaultfontfeatures{Ligatures=TeX,Scale=MatchLowercase}
\fi
% use upquote if available, for straight quotes in verbatim environments
\IfFileExists{upquote.sty}{\usepackage{upquote}}{}
% use microtype if available
\IfFileExists{microtype.sty}{%
\usepackage[]{microtype}
\UseMicrotypeSet[protrusion]{basicmath} % disable protrusion for tt fonts
}{}
\PassOptionsToPackage{hyphens}{url} % url is loaded by hyperref
\usepackage[unicode=true]{hyperref}
\hypersetup{
            pdftitle={Die Bibel kennen: Hebräerbrief},
            pdfborder={0 0 0},
            breaklinks=true}
\urlstyle{same}  % don't use monospace font for urls
\usepackage{natbib}
\bibliographystyle{apalike}
\usepackage{longtable,booktabs}
% Fix footnotes in tables (requires footnote package)
\IfFileExists{footnote.sty}{\usepackage{footnote}\makesavenoteenv{long table}}{}
\usepackage{graphicx,grffile}
\makeatletter
\def\maxwidth{\ifdim\Gin@nat@width>\linewidth\linewidth\else\Gin@nat@width\fi}
\def\maxheight{\ifdim\Gin@nat@height>\textheight\textheight\else\Gin@nat@height\fi}
\makeatother
% Scale images if necessary, so that they will not overflow the page
% margins by default, and it is still possible to overwrite the defaults
% using explicit options in \includegraphics[width, height, ...]{}
\setkeys{Gin}{width=\maxwidth,height=\maxheight,keepaspectratio}
\IfFileExists{parskip.sty}{%
\usepackage{parskip}
}{% else
\setlength{\parindent}{0pt}
\setlength{\parskip}{6pt plus 2pt minus 1pt}
}
\setlength{\emergencystretch}{3em}  % prevent overfull lines
\providecommand{\tightlist}{%
  \setlength{\itemsep}{0pt}\setlength{\parskip}{0pt}}
\setcounter{secnumdepth}{0}
% Redefines (sub)paragraphs to behave more like sections
\ifx\paragraph\undefined\else
\let\oldparagraph\paragraph
\renewcommand{\paragraph}[1]{\oldparagraph{#1}\mbox{}}
\fi
\ifx\subparagraph\undefined\else
\let\oldsubparagraph\subparagraph
\renewcommand{\subparagraph}[1]{\oldsubparagraph{#1}\mbox{}}
\fi

% set default figure placement to htbp
\makeatletter
\def\fps@figure{htbp}
\makeatother

\usepackage[scaled=0.92]{helvet}
\usepackage{booktabs}
\usepackage{amsthm}



\makeatletter
% add Manu
\let\RequirePackage\original@RequirePackage
\let\usepackage\RequirePackage
%
\def\thm@space@setup{%
  \thm@preskip=8pt plus 2pt minus 4pt
  \thm@postskip=\thm@preskip
}
\makeatother

% add Manu
\usepackage{color}
\usepackage{xcolor}
\usepackage{framed}
\setlength{\fboxsep}{.8em}
\usepackage[many]{tcolorbox}
%

% add Manu
\definecolor{aliceblue}{rgb}{0.94, 0.97, 1.0}

\newtcolorbox{blackbox}{
  colback=aliceblue,
  colframe=orange,
  coltext=black,
  boxsep=5pt,
  arc=4pt}

\newenvironment{infobox}
  {
  \begin{itemize}
  \renewcommand{\labelitemi}{
    \raisebox{-.7\height}[0pt][0pt]{
      {\setkeys{Gin}{width=3em,keepaspectratio}
        \includegraphics{images/caution}}
    }
  }
  \setlength{\fboxsep}{1em}
  \begin{blackbox}
  \item
  }
  {
  \end{blackbox}
  \end{itemize}
  }

\newenvironment{definition}{
  \begin{itemize}
    \renewcommand{\labelitemi}{
      \raisebox{-.7\height}[0pt][0pt]{
        {\setkeys{Gin}{width=3em,keepaspectratio}
          \includegraphics{images/caution}
        }
      }
    }
    \setlength{\fboxsep}{1em}
    \begin{blackbox}
      \item
}{
  \end{blackbox}
  \end{itemize}
}

\newenvironment{biblebox}{
  \begin{itemize}
    \renewcommand{\labelitemi}{
      \raisebox{-.7\height}[0pt][0pt]{
        {\setkeys{Gin}{width=3em,keepaspectratio}
          \includegraphics{images/bible}
        }
      }
    }
    \setlength{\fboxsep}{1em}
    \begin{blackbox}
      \item
}{
  \end{blackbox}
  \end{itemize}
}

\newenvironment{quotebox}{
  \begin{itemize}
    \renewcommand{\labelitemi}{
      \raisebox{-.7\height}[0pt][0pt]{
        {\setkeys{Gin}{width=3em,keepaspectratio}
          \includegraphics{images/quote}
        }
      }
    }
    \setlength{\fboxsep}{1em}
    \begin{blackbox}
      \item
}{
  \end{blackbox}
  \end{itemize}
}
%

\title{Die Bibel kennen: Hebräerbrief}
\author{}
\date{\vspace{-2.5em}2020-03-23}

\begin{document}
\maketitle

{
\setcounter{tocdepth}{1}
\tableofcontents
}
\chapter{Studienplan}\label{studienplan}

Eine 12-wöchige praktische Studienreihe über den Hebräerbrief von
Matthew Z. Capps
\href{https://www.thegospelcoalition.org/course/knowing-bible-hebrews}{(Quelle
auf English)}.

Übersetzt aus dem Englischen von eurem Diener.

\begin{itemize}
\tightlist
\item
  \protect\hyperlink{woche01}{Woche 1: Übersicht}
\item
  \protect\hyperlink{woche02}{Woche 2: Einführung: Die Vormachtstellung
  Jesu Christi (Hebr. 1,1-4)}
\item
  \protect\hyperlink{woche03}{Woche 3: Jesus ist den Engelswesen
  überlegen (Hebr. 1,5-2,18)}
\item
  \protect\hyperlink{woche04}{Woche 4: Jesus ist Mose überlegen (Hebr.
  3,1-4,13)}
\item
  Woche 5: Jesus ist der oberste Hohepriester, Teil 1 (Hebr. 4,14-5,10)
\item
  Woche 6: Eine Warnung vor dem Glaubensabfall (Hebr. 5,11-6,20)
\item
  Woche 7: Jesus ist der oberste Hohepriester, Teil 2 (Hebr. 7,1-8,13)
\item
  Woche 8: Jesus ist das überlegene Opfer (Hebr. 9,1-10,18)
\item
  Woche 9: Der Ruf zum Glauben (Hebr. 10,19-11,40)
\item
  Woche 10: Der Ruf zur Ausdauer (Hebr. 12,1-29)
\item
  Woche 11: Letzte Ermahnungen (Hebr. 13,1-25)
\item
  Woche 12: Zusammenfassung und Schlussfolgerung
\end{itemize}

\begin{infobox}
\textbf{Soli Deo Gloria}
\end{infobox}

\begin{center}\rule{0.5\linewidth}{0.5pt}\end{center}

Nur für private Zwecke.

Hebrews: A 12-Week Study © 2015 by Matthew Z. Capps. All rights
reserved.

\url{https://www.thegospelcoalition.org/course/knowing-bible-hebrews/\#week-1-overview}

\hypertarget{woche01}{\chapter{Woche 1: Übersicht}\label{woche01}}

{Custom left tooltip.}

{Custom bottom tooltip.}

{Custom top tooltip.}

{Custom right tooltip.}

Custom bottom tooltip.

\section{Einführung}\label{einfuxfchrung}

Das anonyme Buch der Hebräer ist ein einzigartiger Beitrag zum Kanon der
Schrift. Wie viele andere neutestamentliche Briefe beginnt das
Hebräerbrief ohne Einleitung, schliesst aber mit Segen und Grüssen (Hebr
13,23-24). Der Autor beleuchtet die Form des Hebräerbriefes, indem er
seine Schrift als ``Wort der Ermahnung'' bezeichnet (Hebr 13,22). Das
Hebräische ist mit pastoraler Stimme geschrieben, mit vielen praktischen
Ermahnungen, was viele dazu veranlasst, es als eine einzige Predigt oder
einen einzigen Predigtvortrag zu betrachten, der sich an Bekehrte aus
dem Judentum richtet, die unter dem Druck stehen, zum jüdischen Glauben
zurückzukehren.

Das Hebräerbrief gilt auch als eines der am schönsten geschriebenen und
stilistisch ausgefeiltesten Bücher des Neuen Testaments, ein
literarisches Meisterwerk. Der Autor ist ein Meister der rhetorischen
Debatte und der Überzeugungsarbeit. Er beweist auch seine tiefgreifenden
theologischen Fähigkeiten mit seinem Gebrauch von Bildern, Metaphern,
Anspielungen, alttestamentlichen Analogien und Typologie. In seiner
gesamten Darstellung und Ermahnung webt der Autor einen wunderschönen
Teppich biblischer Theologie mit dem Ziel, die Überlegenheit Jesu
Christi zu verherrlichen.

Das zentrale Motiv der Hebräer ist ``Jesus Christus ist besser'' (die
Worte ``besser'', ``mehr'' und ``grösser'' erscheinen zusammen 25 Mal).
In vielerlei Hinsicht ist die Herrlichkeit Gottes, wie sie sich in Jesus
Christus offenbart, das Gravitationszentrum der Hebräer. Hebräer 1-12
umreisst ein starkes theologisches Argument für die Überlegenheit
Christi über alles Geschaffene und alle Gegenstücke im Alten Testament,
mit besonderem Schwerpunkt auf der Ermutigung des Lesers, in dem
Glauben, der Christus im Zentrum hat, auszuharren. Durch ermutigende
Worte, entschiedene Warnungen und kontrastierende Beispiele ruft der
Autor den Leser oft dazu auf, auf Christus im Gottesdienst zu antworten.

\section{Einordnung in die grössere
Geschichte}\label{einordnung-in-die-gruxf6ssere-geschichte}

Der Hebräerbrief enthält 35 direkte Zitate aus dem Alten Testament und
noch dazu viele Anspielungen und Verweisen auf das alten Testament. Mit
dem alttestamentlichen Hintergrund im Sinn, argumentiert der Autor, dass
Gottes Herrlichkeit und Erlösungsplan schliesslich und am deutlichsten
in Jesus Christus offenbart werden. Die Überlegenheit Jesu zeigt sich
darin, dass er grösser ist als jeder Engel, Priester oder jede
Institution des Alten Bundes. Christus ist das vollständige Sühneopfer
und der letzte Priester. In ihm sehen wir die Erfüllung aller Hoffnungen
und Verheissungen des Alten Testaments, die das lang ersehnte neue
Zeitalter des Bundes einleiten.

\section{Schlüsselvers}\label{schluxfcsselvers}

\begin{quotebox}
{[}Jesus{]} ist die Ausstrahlung seiner Herrlichkeit und der Ausdruck
seines Wesens und trägt alle Dinge durch das Wort seiner Kraft; er hat
sich, nachdem er die Reinigung von unseren Sünden durch sich selbst
vollbracht hat, zur Rechten der Majestät in der Höhe gesetzt.

-- Hebräer 1,3
\end{quotebox}

\section{Datum und historischer
Hintergrund}\label{datum-und-historischer-hintergrund}

Das Hebräische wurde im ersten Jahrhundert, wahrscheinlich vor 70 n.
Chr., geschrieben. Der Autor des Hebräischen nennt sich nicht selbst. Es
gab viele Vermutungen über seine Identität; wie der frühchristliche
Theologe Origenes (ca. 245 n. Chr.) sagte, ``nur Gott weiß'', wer er
ist. Aber wir können sicher sein, dass der Autor mit seinen Zuhörern
vertraut war, denn er sehnte sich danach, mit ihnen wieder vereint zu
werden (Hebr 13,19) und kann ihnen Nachrichten über Timotheus, den
Stellvertreter des Paulus, geben (Hebr 13,23).

Der traditionelle Titel ``An die Hebräer'' spiegelt die alte Vorstellung
wider, dass die ursprüngliche Zuhörerschaft hauptsächlich aus jüdischen
Christen bestand. Man kann mit Sicherheit davon ausgehen, dass die
Zuhörer mit den vielen Zitaten und Anspielungen auf das Alte Testament
vertraut und gut verstanden waren. Sicherlich wandte sich der Autor mit
diesem Brief an bekennende Christen; mehrmals drängt der Autor sie, ihr
Bekenntnis und ihren Glauben aufrechtzuerhalten (Hebr. 3,6.14; Hebr
4,14; Hebr 10,23).

\section{Übersicht}\label{uxfcbersicht}

\begin{enumerate}
\def\labelenumi{\arabic{enumi}.}
\tightlist
\item
  Einführung: Die Vormachtstellung Jesu Christi (Hebr. 1,1-4)
\item
  Jesus ist den Engelswesen überlegen (Hebr. 1,5-2,18)

  \begin{enumerate}
  \def\labelenumii{\roman{enumii}.}
  \tightlist
  \item
    Jesu Status als ewiger Sohn und König (Hebr. 1,5-14)
  \item
    Warnung eins: vor der Vernachlässigung der Errettung (Hebr. 2,1-4)
  \item
    Jesus als der Gründer der Errettung (Hebr. 2,5-18)
  \end{enumerate}
\item
  Jesus ist Mose überlegen (Hebr. 3,1-4,13)

  \begin{enumerate}
  \def\labelenumii{\roman{enumii}.}
  \tightlist
  \item
    Jesus ist größer als Mose (Hebr. 3,1-6)
  \item
    Zweite Warnung: das Scheitern der Exodus-Generation (Hebr. 3,7-19)
  \item
    In Gottes Ruhe eintreten (Hebr. 4,1-13)
  \end{enumerate}
\item
  Jesus ist der oberste Hohepriester, Teil 1 (Hebr. 4,14-5,10)
\item
  Eine Warnung vor dem Glaubensabfall (Hebr. 5,11-6,20)

  \begin{enumerate}
  \def\labelenumii{\roman{enumii}.}
  \tightlist
  \item
    Warnung drei: vor dem Glaubensabfall (Hebr. 5,11-6,12)
  \item
    Die Gewissheit von Gottes Verheißung (Hebr. 6,13-6,20)
  \end{enumerate}
\item
  Jesus ist der oberste Hohepriester, Teil 2 (Hebr. 7,1-8,13)

  \begin{enumerate}
  \def\labelenumii{\roman{enumii}.}
  \tightlist
  \item
    Die Priesterordnung Melchisedeks (Hebr. 7,1-10)
  \item
    Jesus im Vergleich zu Melchisedek (Hebr. 7,11-28)
  \item
    Jesus, ein Priester eines besseren Bundes (Hebr 8,1-13)
  \end{enumerate}
\item
  Jesus ist das überlegene Opfer (Hebr 9,1-10,18)

  \begin{enumerate}
  \def\labelenumii{\roman{enumii}.}
  \tightlist
  \item
    Das irdische Heiligtum (Hebr. 9,1-10)
  \item
    Erlösung durch das Blut Christi (Hebr. 9,11-28)
  \item
    Das Opfer Christi ein für allemal (Hebr 10,1-18)
  \end{enumerate}
\item
  Der Ruf zum Glauben (Hebr. 10,19-11,40)

  \begin{enumerate}
  \def\labelenumii{\roman{enumii}.}
  \tightlist
  \item
    Ermahnung zur Annäherung (Hebr. 10,19-25)
  \item
    Warnung vier: vor dem Zurückschrecken (Hebr. 10,26-39)
  \item
    Durch den Glauben (Hebr. 11,1-40)
  \end{enumerate}
\item
  Der Ruf zur Ausdauer (Hebr. 12,1-29)

  \begin{enumerate}
  \def\labelenumii{\roman{enumii}.}
  \tightlist
  \item
    Jesus, der Gründer und Vervollkommner unseres Glaubens (Hebr 12,1-2)
  \item
    Werde nicht müde (Hebr. 12,3-17)
  \item
    Ein Königreich, das nicht erschüttert werden kann (Hebr 12,18-24)
  \item
    Warnung fünf: vor der Ablehnung des Sprechers (Hebr. 12,25-29)
  \end{enumerate}
\item
  Letzte Ermahnungen (Hebr. 13,1-25)

  \begin{enumerate}
  \def\labelenumii{\roman{enumii}.}
  \tightlist
  \item
    Gott wohlgefällige Opfer (Hebr. 13,1-19)
  \item
    Der Segen (Hebr. 13,20-21)
  \item
    Abschliessende Grüsse (Hebr. 13,22-25)
  \end{enumerate}
\end{enumerate}

\section{Wenn Sie anfangen . . .}\label{wenn-sie-anfangen-.-.-.}

Wie verstehen Sie heute, wie die Hebräer uns helfen, die gesamte
Geschichte der Bibel zu verstehen? Haben Sie eine Vorstellung davon, wie
sich Aspekte des Alten Testaments im Hebräischen erfüllen?

Was ist Ihr gegenwärtiges Verständnis dessen, was die Hebräer zur
christlichen Theologie beitragen? Wie klärt dieses Buch unser
Verständnis der wichtigsten Lehren des christlichen Glaubens?

Gibt es in der hebräischen Sprache eine alttestamentliche Bildsprache,
die Sie besonders verwirrt? Gibt es bestimmte Fragen, von denen Sie
hoffen, dass sie durch diese Studie beantwortet werden können?

\begin{center}\rule{0.5\linewidth}{0.5pt}\end{center}

Nur für private Zwecke. Übersetzt aus dem Englischen von eurem Diener

Hebrews: A 12-Week Study © 2015 by Matthew Z. Capps. All rights
reserved.

\href{https://www.thegospelcoalition.org/course/knowing-bible-hebrews/\#week-1-overview}{source}

\hypertarget{woche02}{\chapter{Woche 2: Einführung: Die Vormachtstellung
Jesu Christi}\label{woche02}}

\begin{biblebox}
\textbf{Lese Hebräer 1,1-4}
\end{biblebox}

\section{Der Ort der Passage}\label{der-ort-der-passage}

Die Anfangsverse der Hebräer stellen Jesus als die letzte und endgültige
Offenbarung Gottes an die Menschheit dar. Der Autor stellt zunächst
fest, dass ``vor langer Zeit, zu vielen Zeiten und auf viele Arten''
Gott zu seinem Volk gesprochen hat (Hebr 1,1). Aber jetzt, in diesen
letzten Tagen, hat Gott endgültig durch Jesus gesprochen - durch seinen
geliebten Sohn, den Schöpfer, Erhalter und Retter der Welt (Hebr 1,2-3)
und ein genaues Bild des Vaters (Hebr 1,3). Die Grösse Jesu wird
schliesslich durch seine Erhöhung zur Rechten Gottes über alle irdischen
und himmlischen Wesen dargestellt (Hebr 1,4).

\section{Das Gesamtbild}\label{das-gesamtbild}

Hebräer 1,1-4 zieht unsere Herzen zur Anbetung im Licht der strahlenden
Majestät und unvergleichlichen Kraft Jesu Christi an.

\section{Reflexion und Diskussion}\label{reflexion-und-diskussion}

Lesen Sie den Abschnitt für diese Studie, Hebräer 1,1-4. Nachdem Sie den
Abschnitt gelesen haben, lesen Sie die folgenden aufgelisteten
Abschnitte erneut und halten Sie Ihre eigenen Antworten auf die Fragen
fest - zuerst in Bezug auf Jesus als göttliche Offenbarung (Hebr.
1,1-2a), dann in Bezug auf Jesu Person, Werk und Status (Hebr. 1,2b-4).

\subsection{1. Jesus als göttliche Offenbarung (Hebr.
1,1-2a)}\label{jesus-als-guxf6ttliche-offenbarung-hebr.-11-2a}

Gott hat gesprochen. In den einleitenden Versen der Hebräer fegt der
Autor über die Spanne von Gottes fortschreitender Offenbarung und landet
auf Jesus Christus als Höhepunkt seiner Kommunikation. Wenn man einige
der wundersamen Weisen betrachtet, auf die Gott im Alten Testament zu
den Patriarchen und Propheten gesprochen hat, was versucht der Autor zu
zeigen, indem er kontrastiert, wie Gott zuvor gesprochen hat und wie
Gott durch seinen Sohn nun endgültig zu seinem Volk gesprochen hat (Hebr
1,1-2)?

In Hebräer 1,1-2 stellt der Autor des Hebräerbriefes die Offenbarung im
Alten Testament der Endgültigkeit der Offenbarung Gottes in Jesus
Christus in vier Bereichen gegenüber. Vergleichen Sie die Epochen der
Offenbarung, die Empfänger der Offenbarung, die Agenten der Offenbarung
und die Art und Weise, wie die Offenbarung ausgedrückt wurde.

Jesus ist die endgültige Offenbarung Gottes in der Geschichte. Diese
Wahrheit impliziert, dass die Offenbarung Gottes im Alten Testament für
diese Epoche ausreichend, aber unvollständig war. Wie wirkt sich die
Offenbarung Jesu ``in diesen letzten Tagen'' darauf aus, wie wir den
vollständigen Kanon der Schrift lesen (Lukas 24,27; Johannes 5,39-40)?

\subsection{2. Jesus: Person, Arbeit und Status (Hebr.
1,2b-4)}\label{jesus-person-arbeit-und-status-hebr.-12b-4}

Viele Gelehrte glauben, dass die Titel ``Sohn'' und ``Erbe'', die in
Hebräer 1,2 auf Jesus angewandt werden, auf Psalm 2,7-8 anspielen, einen
königlichen Krönungspsalm, der an Gottes Versprechen an Davids Erben in
2 Samuel 7,12-16 erinnert. Im alten Israel war es der erstgeborene Sohn,
der das Recht auf das Erbe hatte. Aufgrund seiner königlichen Sohnschaft
ist Jesus der Erbe des Universums, einschließlich der zukünftigen Welt
(Hebr. 2,5-9), was eine Position des Segens und der Ehre ist. Was hebt
der Autor in der einzigartigen Beziehung und Verantwortung Jesu in Bezug
auf das Universum hervor?

In Hebräer 1,2-3 behauptet der Autor, dass die gesamte Schöpfung Gottes
Jesus gehört, weil durch seine Vermittlung alle Dinge entstanden sind
und das Universum durch seine Macht aufrechterhalten wird. Die
Präexistenz, Autorität, Macht und volle Gottheit Jesu sind in seiner
Rolle bei der Erschaffung und Aufrechterhaltung des Universums
offensichtlich (Hebräer 1,10; siehe Johannes 1,3; 1 Korinther 8,6; Kol
1,16). Was vermitteln uns diese Wahrheiten über den Zweck der Schöpfung
und die Herrschaft Jesu über sein Werk?

In Hebräer 1,3 wird Jesus als ``Ausstrahlung der Herrlichkeit Gottes''
beschrieben. In der biblischen Literatur bezieht sich ``Herrlichkeit''
oft auf die leuchtende Manifestation der Person Gottes (siehe Ex. 16,7;
33,18; Jes. 40,5; 60,1, 19). Was die Bedeutung des Wortes ``Glanz''
betrifft, so haben viele festgestellt, dass der Mond Licht reflektiert,
während die Sonne Licht ausstrahlt, weil sie seine Quelle ist. Was sagt
uns das über Jesus als die Ausstrahlung Gottes und unsere Rolle als
Reflektoren der Herrlichkeit Gottes?

In Hebräer 1,3 wird verkündet, dass Jesus ``der genaue Abdruck von
Gottes Natur'' ist. Für die ersten Leser hätte diese Sprache an einen
Eindruck erinnert, der als Bild, wie auf einer Münze, angebracht wurde.
Einfach ausgedrückt: Jesus ist das wahre Abbild Gottes (2 Kor 4,4; Kol
1,15). Wie helfen uns diese Worte, das zu verstehen, was Jesus in
Johannes 14,8-11 lehrte?

Nachdem er sich für die Sünden gereinigt hatte, setzte sich Jesus ``zur
rechten Hand Gottes''. Viele Gelehrte glauben, dass dies eine offene
Anspielung auf Psalm 110,1 ist. Dieser Psalm wird direkt in Hebräer 1,13
zitiert und in Hebräer 8,1; 10,12; 12,2 angedeutet. Was wird mitgeteilt,
wenn der Autor schreibt, dass Christus ``zur Rechten Gottes gesessen
hat''?

Jesus wird nicht nur den Propheten, sondern auch den Engeln wegen seines
vorzüglicheren geerbten Namens für überlegen erklärt (Hebr. 1,4). Der
Autor scheint 2. Samuel 7 in Bezug auf die Ehre, die Jesus als
Davidischer Erbe zuteil wurde, zu wiederholen. Welche Bedeutung hat es,
Jesus in der Position des königlichen Erben von den Engeln zu
unterscheiden?

\begin{definition}
\textbf{Engel -- Definition}

Ein übernatürlicher Gesandter Gottes, oft gesandt, um seinen Willen zu
erfüllen oder den Menschen bei der Erfüllung seines Willens zu helfen.
Obwohl Engel mächtiger sind als Menschen und oft Ehrfurcht einflößen,
sollen sie nicht angebetet werden (Kol. 2,18; Offb. 22,8-9). In der
Bibel gibt es jedoch verschiedene Erscheinungen eines ``Engels des
Herrn'', der offenbar eine physische Manifestation Gottes selbst ist.
\end{definition}

Lesen Sie die folgenden drei Abschnitte zu den Themen
Evangeliums-Blicke, Ganz-Bibel-Verbindungen und Theologische
Sondierungen. Nehmen Sie sich dann Zeit, über die persönlichen
Auswirkungen nachzudenken, die diese Abschnitte für Ihren Weg mit dem
Herrn haben können.

\section{Einblicke in das Evangelium}\label{einblicke-in-das-evangelium}

\textbf{REINIGUNG DER SÜNDE.} Die Bibel macht deutlich, dass die Sünde
und ihre Verderbnis zerstörerische Auswirkungen auf die Menschheit und
die gesamte Schöpfung haben. Die Notwendigkeit der Reinigung von der
Sünde ist Teil der übergreifenden Handlung der Bibel. Das kosmische
Ausmaß der Sünde bildet die Grundlage für die kosmische Erlösung durch
den Sühnetod Jesu. Im Alten Testament wurden Sühneopfer für das Volk
Gottes eingeführt, um die Strafe für seine Sünde zu vermitteln und die
Reinigung durch Blut zu erlangen (Levitikus 16). Sühneopfer reinigten
auch die Gegenstände des irdischen Tempels, der dem Kosmos nachempfunden
ist, der Tempelbehausung Gottes. Im Hebräerbrief sehen wir, dass der Tod
Jesu die notwendige Reinigung von Sünden und die Reinigung des
menschlichen Gewissens vor Gott erlangte (Hebr. 1,3; 9,14). Das Opfer
Jesu erstreckte sich auch auf die Reinigung ``himmlischer Dinge'' (Hebr.
9,23) und ist daher kosmisch. Das Herzstück des Evangeliums ist die gute
Nachricht, dass das Blut Jesu für die Sünde gesühnt hat und Auswirkungen
auf den gesamten Kosmos hat.

\textbf{DIE VOLLSTÄNDIGE ERRETTUNG.} Kurz bevor Jesus seinen letzten
Atemzug am Kreuz tat, verkündete er: ``Es ist vollbracht'' (Johannes
19,30). Das Werk, zu dessen Vollendung der Vater ihn gesandt hatte, war
vollendet - sein vollkommenes Opfer für unsere Sünde war vollendet
(Hebr. 1,3; 9,11-12, 25-28). Der Verfasser des Hebräerbriefes weist
darauf hin, dass sich Jesus, nachdem sein Werk am Kreuz und in der
Auferstehung vollendet war, zur rechten Hand Gottes setzte und damit die
Endgültigkeit seines Werkes und seines Status unterstrich. Im Gegensatz
zu den levitischen Priestern, die Jahr für Jahr unvollkommene Opfer
brachten, um die Sünden zu verdecken, brachte Jesus das perfekte
Einwegopfer dar, das die Sünde auslöschte - und dann nahm er seinen
Platz ein, um für immer zu regieren (Hebr 10,11-12).

\begin{definition}
\textbf{Sünde -- Definition}

Jede Verletzung oder Nichtbeachtung der Gebote Gottes oder der Wunsch,
dies zu tun.
\end{definition}

\section{Ganz-Bibel-Verbindungen}\label{ganz-bibel-verbindungen}

\textbf{OFFENBARUNG UND ERLÖSUNG.} Offenbarung in menschlicher Sprache
ist wesentlich für die Vermittlung von Gottes Erlösungsplan durch Jesus
Christus. Ohne verbale Offenbarung kann die Menschheit keinen Zugang zu
der guten Nachricht von Gottes Erlösung haben. Die Offenbarung im Alten
Testament steht jedoch nicht für sich allein, sie ist unvollständig ohne
ihren Abschluss und ihre Erfüllung in Jesus Christus. ``In diesen
letzten Tagen'' kommen wir zu der Erkenntnis, dass die Personen und
Institutionen des Alten Testaments auf die Person und das Werk Jesu
Christi hinweisen, in dem wir die Erlösung finden (Hebr 1,1-2).

\textbf{DEN VOLLKOMMENEN SOHN.} In der Bibel wird Sohnschaft mit
Familienähnlichkeit, Abstammung und Erbe in Verbindung gebracht. Nicht
nur Gottes erster ``Sohn'' Adam wurde nach seinem Bilde geschaffen, Adam
gebar auch Söhne nach seinem eigenen Bilde und damit nach dem Bilde
Gottes (1. Mose 1,28; 5,1-3). Gott bezeichnet Israel später als seinen
gemeinsamen ``Sohn'' (Ex 4,22-23; Ps 2,7; Hos 11,1) und seinen
``Erstgeborenen'' (Deut 33,17; Ps 2,7; Jer 31,9; Esra 6,58). Sowohl Adam
als auch Israel waren nicht das, was der Vater wollte. Beide taten
nicht, was Gott von ihnen als ``Söhne'' verlangt hatte. Ihr Ungehorsam
steht in krassem Gegensatz zu dem makellosen Gehorsam Jesu, des
göttlich-menschlichen Sohnes, der den Vater vollkommen nachbildet und
ihm Ehre bringt (Spr 10,1; 15,20; 23,15).

\textbf{DIE HERRLICHKEIT GOTTES.} In der biblischen Sprache ist die
Herrlichkeit Gottes ein Abbild seiner Vollkommenheit, Schönheit und
Größe. Mose 1,27 wird uns gesagt, dass Adam als Gottes Ebenbild
geschaffen wurde. Als Gottes Bildträger wurde Adam - zusammen mit dem
Rest der Menschheit - geschaffen, um Gottes Herrlichkeit
widerzuspiegeln. Aber die Sünde zerstörte die reine Widerspiegelung der
Herrlichkeit Gottes in Adam und seinen Kindern. Von der Sünde
unvermindert wird Jesus als der zweite und letzte Adam erklärt, der
Gottes Ebenbild vollständig repräsentiert und seine Herrlichkeit
makellos ausstrahlt (Römer 5,12-21). Doch im Gegensatz zu Adam ist Jesus
die genaue Prägung Gottes und in der Substanz identisch mit Gott (Hebr.
1,3).

\section{Theologische Sondierungen}\label{theologische-sondierungen}

\textbf{SCHÖPFUNG.} Die biblische Geschichte beginnt mit einer
majestätischen Beschreibung, wie Gott Himmel und Erde zu seiner
Behausung schuf und wie sie von seinen Geschöpfen bewohnt wurde. In
Hebräer 1,2-3 wird uns gesagt, dass Jesus nicht nur das Instrument des
ursprünglichen Schöpfungsaktes in Genesis 1 ist, sondern auch eng in die
fortgesetzte Sorge für die Schöpfung eingebunden ist (vgl. die
Beschreibung Jesu in Johannes 1,1-18). Jesus war der Akteur, in dem und
durch den das gesamte Universum von Raum und Zeit entstand. Als das
Abbild des menschgewordenen Gottes ist Jesus der Berührungspunkt
zwischen dem Schöpfer und seinem Universum. Er ist der Bezugsrahmen für
den ursprünglichen Zweck und für die Erneuerung der Schöpfung Gottes
nach dem Fall (Kol. 1,15-23).

\textbf{VORSEHUNG.} Die Vorsehungslehre lehrt, dass Gott die von ihm
geschaffene Welt erhält und sie zu den von ihm festgelegten Zielen
führt. In Hebräer 1,3 sehen wir, dass Gott persönlich an seiner
Schöpfung beteiligt ist, um sie zu erhalten und zu bewahren. Seine
providentielle Herrschaft erstreckt sich über alle Dinge im Universum;
die gesamte Schöpfung wird durch das mächtige Wort Jesu erhalten und
weitergeführt. Darüber hinaus kam Jesus, um die endgültige und
vollständige Reinigung für die Sünden zu gewährleisten.

\section{Persönliche Implikationen}\label{persuxf6nliche-implikationen}

Nehmen Sie sich Zeit, über die Auswirkungen von Hebräer 1,1-4 auf Ihr
eigenes Leben heute nachzudenken. Beachten Sie die persönlichen
Implikationen für Ihren Weg mit dem Herrn im Licht der (1)
Evangeliums-Blicke, (2) Gesamtbiblischen Zusammenhänge, (3)
Theologischen Klänge und (4) dieses Abschnitts als Ganzes.

\begin{enumerate}
\def\labelenumi{\arabic{enumi}.}
\tightlist
\item
  Einblicke in das Evangelium
\item
  Ganz-Bibel-Verbindungen
\item
  Theologische Sondierungen
\item
  Hebräer 1,1-4
\end{enumerate}

\section{Wenn Sie diese Einheit abschliessen . .
.}\label{wenn-sie-diese-einheit-abschliessen-.-.-.}

Nehmen Sie sich jetzt einen Moment Zeit, um den Segen und die Hilfe des
Herrn zu erbitten, während Sie mit dem Studium der Hebräer fortfahren.
Und nehmen Sie sich auch einen Moment Zeit, um durch diese
Studieneinheit zurückzublicken, über einige Dinge nachzudenken, die der
Herr Sie vielleicht lehrt, und Dinge zu notieren, die Sie in Zukunft
noch einmal überprüfen sollten.

\begin{center}\rule{0.5\linewidth}{0.5pt}\end{center}

Nur für private Zwecke. Übersetzt aus dem Englischen von eurem Diener

Hebrews: A 12-Week Study © 2015 by Matthew Z. Capps. All rights
reserved.

\href{https://www.thegospelcoalition.org/course/knowing-bible-hebrews/\#week-2-introduction-the-supremacy-of-jesus-christ-heb-11-4}{source}

\hypertarget{woche03}{\chapter{Woche 3: Jesus ist den Engelswesen
überlegen}\label{woche03}}

\begin{biblebox}
\textbf{Lese Hebräer 1,5-2,18}
\end{biblebox}

\section{Der Ort der Passage}\label{der-ort-der-passage-1}

Durch eine Kette von alttestamentlichen Passagen wird Jesus als der
einzige Sohn Gottes den Engeln überlegen gezeigt (Hebr. 1,5-14).
Aufgrund seines Status als Sohn Gottes genießt Jesus eine einzigartige
Beziehung zum Vater in Position, Wesen und Autorität. Jesus wird auch
als der einzigartige Menschensohn dargestellt, was ihn als das wahre
Opfer für die Sünde etabliert und ihn als den mitfühlenden Hohenpriester
einführt (Hebr. 2,1-18). Dieser Abschnitt enthält die erste von fünf
Warnungen im Buch der Hebräer; hier werden wir gewarnt, auf die durch
Jesus Christus bereitgestellte Rettung zu vertrauen.

\section{Das Gesamtbild}\label{das-gesamtbild-1}

Hebräer 1,5-2,18 erhebt Jesus über die Engel als den einzigartigen Sohn
Gottes, der die Sünden der Menschheit versühnte und nun als Mittler und
Hoherpriester zwischen Gott und der Menschheit dient.

\begin{definition}
\textbf{Sühne -- Definition}

Die Beschwichtigung des Zorns durch die Darbringung einer Gabe oder
eines Opfers. Jesus versöhnte die Sünden der Menschheit durch sein
Leiden und seinen Tod (Römer 3,25; Hebr 2,17; 1. Johannes 2,2; 4,10).
\end{definition}

\section{Reflexion und Diskussion}\label{reflexion-und-diskussion-1}

Lesen Sie den gesamten Abschnitt für diese Studie, Hebräer 1:5-2:18.
Nachdem Sie den gesamten Abschnitt gelesen haben, beginnen Sie damit,
die aufgelisteten Teile erneut durchzulesen und Ihre eigenen Antworten
auf die folgenden Fragen aufzuzeichnen. (Weitere
Hintergrundinformationen finden Sie in der ESV-Studienbibel, Seiten
2361-2364, die online unter www.esvbible.org verfügbar ist).

\subsection{1. Der Status Jesu als ewiger Sohn und König (Hebr.
1,5-14)}\label{der-status-jesu-als-ewiger-sohn-und-kuxf6nig-hebr.-15-14}

Der Autor unterstützt die Überlegenheit Jesu gegenüber den Engeln, indem
er verschiedene Texte des Alten Testaments aneinander reiht, um seine
Argumentation zu untermauern. Jesus wird durch die Art seiner Beziehung
zum Vater (Hebr. 1,5), seine Stellung über die Engel (Hebr. 1,6-7) und
seine königliche Autorität (Hebr. 1,8-9) als überlegen dargestellt. Wie
stärkt die Anwendung alttestamentlicher Texte durch den Autor auf Jesu
Sohnschaft (Ps. 2,7; 2 Sam. 7,14), seine Stellung (Ps. 97,7; 104,4) und
seine Autorität (Ps. 45,6-7) unseren Glauben an Jesus heute?

Mit dem Zitat von Psalm 102,25-27 in Hebräer 1,10-12 betont der Autor
die Rolle Jesu in der Schöpfung und seine ewige Natur. Wie stärken Jesu
Rolle in der Schöpfung und seine ewige Herrschaft über die Schöpfung
unsere Treue zu ihm?

In Hebräer 1,13 wendet der Autor Psalm 110,1 auf Jesus an und zeigt ihn
als erhaben zur rechten Hand Gottes, eine Position des Vorrechts und der
Macht (siehe Hebräer 1,3). Dieser spezielle Psalm bezieht sich auf die
Inthronisierung des Königs und den Sieg über alle seine Feinde. Was
vermittelt die Sitzhaltung Jesu über die Absichten Gottes im Leben, Tod
und in der Auferstehung Christi?

Nach Hebräer 1,14 sind Engel dienende Geister, die die Aufgabe haben,
denen zu dienen, die die Erlösung erben sollen. Wie spricht die Rolle
der Engel zur Ehre des Evangeliums und zur Autorität von Jesus Christus?

\subsection{2. Warnung Eins: Vor der Vernachlässigung der Errettung
(Hebr.
2,1-4)}\label{warnung-eins-vor-der-vernachluxe4ssigung-der-errettung-hebr.-21-4}

Der Autor untermauert die Verlässlichkeit des mosaischen Gesetzes
dadurch, dass es von den Engeln gegeben wurde (Hebr. 2,2; siehe auch
Deut. 33,2). Wenn das mosaische Gesetz mit einer Vergeltung für
Ungehorsam kam, wie viel mehr werden dann diejenigen belohnt, die das
von Gott dem Vater und Gott dem Sohn bezeugte und durch Zeichen, Wunder
und Wundertaten bestätigte Heil ablehnen (Hebr. 2,3-4)? Warum sollten
wir außerdem darauf achten, dass wir nicht vom Evangelium der Gnade
abweichen und seine Anwendung auf unser Leben bewusst vernachlässigen
(Hebr. 2,1)?

\subsection{3. Jesus als der Gründer des Heils (Hebr.
2,5-18)}\label{jesus-als-der-gruxfcnder-des-heils-hebr.-25-18}

In Hebräer 2,5-9 heißt es, dass die gegenwärtige und die zukünftige Welt
Jesus Christus unterworfen sind. Aber zum gegenwärtigen Zeitpunkt sehen
die Gläubigen die höchste Herrschaft Jesu über den Kosmos nicht klar.
Darüber hinaus sind die Menschen angesichts des Sündenfalls und des
Versagens des Menschen, den Schöpfungsauftrag aufrechtzuerhalten,
vorübergehend in ihrem Status und ihrer Autorität niedriger als die
Engel (1. Mose 1,28). Wie erfüllt Jesus als der wahre Vertreter der
Menschheit Gottes Gebot, alles in der Schöpfung unter die Herrschaft zu
stellen (Hebr 2,8)?

In Hebräer 2,10-13 zitiert der Autor Psalm 22,22 und Jesaja 8,17b-18, um
zu zeigen, dass die Nachfolger des einen einzigartigen Sohnes Gottes nun
auch ``Söhne'' genannt werden, denn sie werden durch Jesu vollkommenes
Leben und Opfer in die neu erlöste Menschheitsfamilie aufgenommen.
Welche Vorteile hat es, ein Sohn Gottes oder ein Bruder Jesu Christi zu
sein (siehe Gal 4,1-7)?

In Hebräer 2,14-18 sehen wir die Solidarität Jesu mit der Menschheit,
indem er ``Fleisch und Blut'' annimmt. Doch im Gegensatz zu dem, was
jeder andere Mensch hätte tun können, stürmte Jesus die Tore des Todes,
besiegte den Bösen und befreite uns aus der Sklaverei bis zum Tod (Hebr.
2,14-16). Wie also gibt Hebräer 2,17-18 den Gläubigen inmitten
geistlicher Gebrechen Hoffnung?

\begin{definition}
\textbf{Gnade -- Definition}

Unverdiente Gunst, insbesondere die kostenlose Gabe der Erlösung, die
Gott den Gläubigen durch den Glauben an Jesus Christus schenkt.
\end{definition}

Lesen Sie die folgenden drei Abschnitte zu den Themen
Evangeliums-Blicke, Ganz-Bibel-Verbindungen und Theologische
Sondierungen. Nehmen Sie sich dann Zeit, über die persönlichen
Auswirkungen nachzudenken, die diese Abschnitte für Ihren Weg mit dem
Herrn haben können.

\section{Einblicke in das
Evangelium}\label{einblicke-in-das-evangelium-1}

\textbf{ZUVERLÄSSIGES EVANGELIUM.} Die Gefahr, das Evangelium Jesu
Christi zu vernachlässigen, wird durch seine Überlegenheit gegenüber der
vorherigen Offenbarung, die durch Propheten und Engel kam, noch erhöht
(Hebr. 1,2; 2,2). Verlässlichen Augenzeugen zufolge verkündete Jesus
selbst zuerst seine gute Nachricht (Hebr. 2,3). Darüber hinaus
bestätigte Gott diese von Jesus verkündete gute Nachricht mit Zeichen,
Wundern, Wundertaten und Gaben des Heiligen Geistes (Apg 2,22; 2. Kor
12,12; Hebr 2,4). Das Evangelium ist nicht nur eine gute Nachricht; es
ist eine verlässliche gute Nachricht - eine gute Nachricht, auf die es
sich lohnt, unser Leben zu riskieren.

\textbf{DER GUTE KÖNIG.} Im ersten Jahrhundert wurde das ``Evangelium''
(``gute Nachricht'') regelmäßig verwendet, um die Geburt, die
Verkündigung, die Thronbesteigung oder den Sieg eines großen Königs zu
bezeichnen. Nach Hebräer 1,8 ist Jesus der ewige messianische König,
dessen Herrschaft nie enden wird. Im Gegensatz zu den davidischen
Königen der Vergangenheit wird seine Herrschaft nicht durch
Zerbrechlichkeit behindert. Die Herrschaft Jesu ist von vollkommener
Gerechtigkeit und Rechtschaffenheit geprägt. Die königliche Herrschaft
Jesu ist auch insofern gut, als er sein Volk befreit, heiligt und für es
sorgt (Hebr. 1,8; 2,6; 9-18). Schließlich wird Jesus, anders als jeder
andere König in der Geschichte, für immer und ewig als König der Könige
regieren (Offb 5,9-14).

\textbf{TRANK TIEF AUS DEM TOD.} Als vollkommenes Opfer kostete Jesus
den Tod für alle, die glauben (Hebr. 2,9). Mehr noch, Jesus trank den
Kelch des Zornes Gottes bis zum bitteren Abschaum, um den Zorn Gottes
für die Gläubigen zu verzehren und so den Todesgriff des Teufels zu
zerstören (Hebr. 2,14-15, 17). In Christus gibt es keine Angst im Tod,
sondern nur Hoffnung im Leben. Am Kreuz hat Jesus die Macht des Todes
beseitigt. In seiner mächtigen Auferstehung besiegelte Jesus die
Verheißung eines neuen ewigen Lebens. Ganz-Bibel-Verbindungen

\textbf{VERWALTETER KÖNIG.} In Hebräer 1,5b greift der Autor die
Verkündigung König Davids über seinen Bundeserben auf, den Gott als
seinen eigenen Sohn einsetzen wird (2 Sam 7,14; 1 Chron 17,13). Im
Zusammenhang mit 2 Samuel ist dies sicherlich Salomo. Während diese
Worte in Salomo nie verwirklicht wurden, haben sie ihre Erfüllung in
Jesus gefunden, der der wahre und größere Davidskönig ist, dessen
Königreich für immer errichtet ist (2. Sam. 7,16). Daher findet die
königliche Herrschaft, die die israelitischen Könige vorausgesagt haben
- die aber nicht als Vertreter Gottes auf Erden aufrecht erhalten werden
konnte - ihre Erfüllung in Jesus (Hebr. 2,8-9; Offb. 3,21). Die Könige
des Alten Testaments sind die Schatten; Jesus ist die Substanz. Jesus
ist der ewige König der Könige, dessen vollkommene Herrschaft und
Gerechtigkeit niemals enden wird.

\textbf{DER GEISTLICHE EXODUS.} Der Autor des Hebräerbriefes bezeichnet
Jesus nicht nur als den Gründer des Heils, sondern auch als denjenigen,
der das Volk des neuen Bundes in die Herrlichkeit führt (Hebr. 2,10).
Einige Bibelversionen übersetzen ``Gründer'' mit ``Pionier''. Die
Symbolik von Jesus, der unsere Rettung vorbereitet hat, blickt im Alten
Testament auf diejenigen zurück, die die Israeliten durch die Wüste und
auch in die Schlacht geführt haben (Num. 10,4; 13,2-3; Jud. 5,15; 9,44;
11,6-11; 1 Chron. 5,24; 26,26; 2. Chron. 23,14; Neh. 2,9). Als Vorläufer
und Vertreter seines Volkes ist Jesus in die Gegenwart Gottes
eingetreten, um ihnen den Eintritt in die Gegenwart Gottes zu sichern;
er ist der Weg geworden, auf dem sie in die verheißene Ruhe Gottes
eintreten können. Jesus ist der wahre und größere Mose, der sein ganzes
Volk in Gottes Ruhe erlöst.

\textbf{BESCHAFFUNG.} Das Wort ``Versöhnung'' vermittelt den Sinn eines
Sühneopfers, das den Zorn Gottes gegen die Sünde befriedigt (Römer 3,25;
Hebr 2,17). In der gesamten Erlösungsgeschichte wird Gottes gerechter
Zorn als Notwendigkeit gezeigt, besänftigt zu werden, bevor die Sünde
des Volkes Gottes vergeben werden kann. Die endgültige Versöhnung durch
Christus, die ein für allemal den vollen Lohn für die Sünde bringt, wird
im Alten Testament mehrmals angedeutet (Ex 32,11-14; Num 25,8; Josh
7,25-26). Wo die alttestamentlichen Opfer versagten, gelang es Jesus,
Gottes Zorn ein für allemal gegen die Sünde seines Volkes zu versöhnen.

\begin{definition}
\textbf{Erlösung -- Definition}

Befreiung von den ewigen Folgen der Sünde. Der Tod und die Auferstehung
Jesu erkauften den Gläubigen ewige Erlösung (Römer 1,16).
\end{definition}

\section{Theologische Sondierungen}\label{theologische-sondierungen-1}

\textbf{ANGELS.} Nach der Heiligen Schrift sind Engel majestätische,
erschaffene Wesen, die in erster Linie als Anbeter und Boten Gottes
fungieren, seinen Willen offenbaren und Schlüsselereignisse in der
gesamten Erlösungsgeschichte ankündigen (Dan 9,20-27; Lukas 1,11-20; Apg
7,38; Hebr 2,2). Engel dienen auch dem Volk Gottes (1. Könige 19,5-7;
Ps. 91,11; Hebr. 1,14). Die Pracht der Engel dient als Bezugspunkt, von
dem aus der Autor der Hebräer von der viel höheren Stellung des
erhabenen Sohnes sprechen kann (Hebr. 1,5-13). Die Anbetung Christi
durch die Engel, wenn er in den Himmel kommt, festigt seine überragende
Stellung (Hebr. 1,6). Obwohl Engel schöne und mächtige Wesen sind,
übertrifft die Schönheit und Macht Christi die ihre exponentiell.

\textbf{MENSCHLICHKEIT JESU.} Während seiner Zeit auf der Erde wurde
Jesus als der inkarnierte Gottmensch niedriger als die Engel gemacht,
indem er voll an menschlichem Fleisch und Blut teilhatte (Hebr. 2,9.14).
Die Menschlichkeit Christi war notwendig, damit er Versuchungen und
Leiden ertragen und so als das einzige wahre Opfer für die Sünde dienen
konnte. Wenn Jesus nicht in jeder Hinsicht (außer der Sünde) menschlich
wurde und die ganze Bandbreite der Versuchung erfuhr, konnte er nicht
als mitfühlender Hohepriester dienen, der die geistlichen Gebrechen der
Menschen kennt.

\textbf{WUNDER.} In der Bibel sind Wunder Gottes nicht-normative
Machttaten, durch die er von sich selbst Zeugnis ablegt und seine Boten
und seine Botschaft beglaubigt (Johannes 2,11; 3,2; Apostelgeschichte
2,22). In der frühen Kirche vollbrachten die Apostel und andere
ebenfalls Wunder, um die Gültigkeit der von ihnen verkündeten
Evangeliumsbotschaft zu bestätigen (Apg 2,43; 3,6-10; 4,30; 8,6-8,13;
9,40-42). In Hebräer 2,3-4 lesen wir, dass Gott die Botschaft der
Erlösung selbst durch Zeichen und Wunder bestätigt hat.

\section{Persönliche
Implikationen}\label{persuxf6nliche-implikationen-1}

Nehmen Sie sich Zeit, über die Auswirkungen von Hebräer 1,5-2,18 auf Ihr
eigenes Leben heute nachzudenken. Beachten Sie die persönlichen
Implikationen für Ihren Weg mit dem Herrn im Licht der (1)
Evangeliums-Blicke, (2) Gesamtbiblischen Zusammenhänge, (3)
Theologischen Klänge und (4) dieses Abschnitts als Ganzes.

\begin{enumerate}
\def\labelenumi{\arabic{enumi}.}
\tightlist
\item
  Einblicke in das Evangelium
\item
  Ganz-Bibel-Verbindungen
\item
  Theologische Sondierungen
\item
  Hebräer 1,5-2,18
\end{enumerate}

\section{Wenn Sie diese Einheit abschliessen . .
.}\label{wenn-sie-diese-einheit-abschliessen-.-.-.-1}

Nehmen Sie sich jetzt einen Moment Zeit, um den Segen und die Hilfe des
Herrn zu erbitten, während Sie mit dem Studium der Hebräer fortfahren.
Nehmen Sie sich auch einen Moment Zeit, um diese Studieneinheit noch
einmal durchzugehen, über einige Dinge nachzudenken, die der Herr Sie
vielleicht lehrt, und sich Dinge zu merken, die Sie in Zukunft noch
einmal überprüfen sollten.

\begin{center}\rule{0.5\linewidth}{0.5pt}\end{center}

Nur für private Zwecke. Übersetzt aus dem Englischen von eurem Diener

Hebrews: A 12-Week Study © 2015 by Matthew Z. Capps. All rights
reserved.

\href{https://www.thegospelcoalition.org/course/knowing-bible-hebrews/\#week-3-jesus-is-superior-to-angelic-beings-heb-15-218}{source}

\hypertarget{woche04}{\chapter{Woche 4: Jesus ist Mose
überlegen}\label{woche04}}

\begin{biblebox}
\textbf{Lese Hebräer 3,1-4,13}
\end{biblebox}

\section{Der Ort der Passage}\label{der-ort-der-passage-2}

In Hebräer 3,1-6 wird gezeigt, dass Jesus Moses, einem der treuesten
Knechte Gottes, überlegen ist. Jesus ist der höchsten Ehre würdig, weil
er der treue Hohepriester und Sohn Gottes ist. Der Autor ermahnt die
Christen, mit Treue und Ausdauer auf Gottes Erlösungswerk zu antworten
(Hebr. 3,7-4,13). In diesem Abschnitt erhalten die Leser ihre zweite
Warnung, nämlich die, im Gegensatz zu den Menschen der
Exodus-Generation, die ihre Herzen verhärtet haben, im Glauben
auszuharren (Hebr 3,7-18).

\section{Das grosse Bild}\label{das-grosse-bild}

Hebräer 3,1-4,13 zeigt Jesus Christus als den Apostel und vollkommenen
Sohn, der gesandt wurde, um als treuer Hohepriester zu dienen.

\section{Reflexion und Diskussion}\label{reflexion-und-diskussion-2}

Lesen Sie den gesamten Abschnitt für diese Studie, Hebräer 3,1-4,13.
Nachdem Sie den gesamten Abschnitt gelesen haben, beginnen Sie damit,
die aufgelisteten Teile erneut durchzulesen und Ihre eigenen Antworten
auf die folgenden Fragen aufzuzeichnen.

\subsection{1. Jesus ist grösser als Moses (Hebr.
3,1-6)}\label{jesus-ist-gruxf6sser-als-moses-hebr.-31-6}

In Hebräer 3,1-2 bekennt sich Jesus als der Apostel, der gesandt wurde,
um Befreiung zu bringen, und als der Hohepriester, der die Sünden des
Volkes Gottes sühnen soll. Was bedeutet es, als heilige Mitglieder der
Familie Gottes ``an einer himmlischen Berufung teilzuhaben'' (Hebräer
3,1)? Wie kann die Treue Jesu uns ermutigen, im Vertrauen und in der
Hoffnung an unserer himmlischen Berufung festzuhalten (Hebr. 3,6)?

Mose war ein treuer Knecht im Haus Gottes trotz der Treulosigkeit der
Israeliten (Num 12,7). Aber Jesus ist wegen seiner vollkommenen Treue
als der Sohn, der nicht nur dem Haus Gottes vorsteht, sondern es auch
gebaut hat, einer grösseren Ehre würdig als Mose (1. Chron. 17,14-17;
Hebr. 3,1-6). Welche Bedeutung hatte für die ursprüngliche Zuhörerschaft
- die mit dem Alten Testament vertraut war - der Vergleich des Autors
mit Jesus und Mose? Was können wir von der Treue des Mose als Mitglied
des Hauses Gottes lernen?

\subsection{2. Warnung Zwei: Das Versagen der Exodus-Generation (Hebr.
3,7-19)}\label{warnung-zwei-das-versagen-der-exodus-generation-hebr.-37-19}

Nachdem der Autor den Kontrast zwischen der Treue von Moses und Jesus
untersucht hat, wendet er sich nun den Antworten ihrer Anhänger zu,
indem er Psalm 95,7-11 zitiert. Der Autor warnt vor dem Unglauben eines
sündigen, verhärteten Herzens, das einen zum Abfall veranlasst (Hebr.
3,7-12). Wie tragen diese starken Worte dazu bei, uns davon abzuhalten,
in der Rebellion zu leben? Wie erleuchtet uns dieser Abschnitt über den
Unterschied zwischen echtem und falschem Glauben?

In Hebräer 3,13-14 erklärt der Autor, wie gegenseitige Ermutigung und
Rechenschaftspflicht das Gegenmittel gegen die Zersetzung des Unglaubens
sein können. Wie hält die Befolgung der Ermahnung zu gegenseitiger
Verpflichtung im Teilen Christi Christen dazu an, bis zum Ende
durchzuhalten?

In Hebräer 3,15 kehrt der Autor zu Psalm 95,7-8 zurück und erinnert die
Leser an die Misserfolge der Exodusgeneration (Hebräer 3,16-19; siehe Ex
17,1-7; Num 14,20-38). Wie weist der Unglaube der Exodus-Generation und
ihr Versagen, in Gottes Ruhe zu gehen, sowohl auf die Warnung vor dem
Abfall als auch auf die Hoffnung hin, die wir in Jesus Christus haben?

\subsection{3. In Gottes Ruhe eintreten (Hebr.
4,1-13)}\label{in-gottes-ruhe-eintreten-hebr.-41-13}

Gottes Befreiung der Israeliten und ihre Hoffnung, das verheissene Land
zu betreten, war ein Vorzeichen für das Werk Jesu. Für die Kirche ist
die verheissene Ruhe Gottes in der Auferstehung Christi eingeweiht
worden und wartet auf die Vollendung bei seinem zweiten Kommen; aber der
Autor fürchtet, dass einige nicht in diese Ruhe eintreten werden (Hebr
4,1). Wie verhält sich der Glaube nach Hebräer 4,1-3 dazu, in die
verheissene Ruhe Gottes einzutreten?

In Hebräer 4,4 spricht der Autor von der Einkehr in Gottes Ruhe in Bezug
auf den siebten Schöpfungstag (1. Mose 2,2) und stellt fest, dass die
Israeliten es letztlich versäumt haben, in die Ruhe Gottes einzutreten
(Hebräer 4,5-9). Zusammengenommen ist Gottes Sabbatruhe, die in Genesis
2 begann, immer noch offen und kann betreten werden. Was bedeutet es, in
seine Werke einzutreten und von ihnen zu ruhen, wie Gott von seinen
Werken (Hebr. 4,10)?

Nachdem der Autor die Leser aufgerufen hat, sich darum zu bemühen, in
Beharrlichkeit in Gottes Ruhe einzutreten, erinnert er sie daran, dass
treuloser Ungehorsam nicht unbemerkt bleiben wird (Hebr. 4,11-13). Wenn
wir wissen, dass Gottes Wort wie ein Mittel wirkt, das die innersten
Gedanken und Absichten offenbart, wie können wir dann sein Wort
benutzen, um Ungehorsam und Unglauben zu bekämpfen?

\begin{definition}
\textbf{Sabbat -- Definition}

Samstag, der siebte Tag der Woche, der jüdische Tag der Anbetung und der
Ruhe (1. Mose 2,2-3; Ex 31,13-17). Christen treffen sich am Sonntag, dem
Tag der Auferstehung Christi (Apg 20,7), zum Gottesdienst und betrachten
den Sonntag statt des Samstags als ihren wöchentlichen Ruhetag. Und doch
freuen sich die Gläubigen auf eine ewige Sabbatruhe (Hebr 4,1-13).
\end{definition}

\begin{definition}
\textbf{Werke -- Definition}

Handlungen und Einstellungen, entweder gut oder schlecht. Wahrer Glaube
an Christus wird unweigerlich gute Werke hervorbringen, die Gott
wohlgefällig sind. Gute Werke können niemals die Grundlage oder das
Mittel zur Erlösung sein, die allein aus Gnade durch den Glauben
erfolgt.
\end{definition}

Lesen Sie die folgenden drei Abschnitte zu den Themen
Evangeliums-Blicke, Ganz-Bibel-Verbindungen und Theologische
Sondierungen. Nehmen Sie sich dann Zeit, über die persönlichen
Auswirkungen nachzudenken, die diese Abschnitte für Ihren Weg mit dem
Herrn haben können.

\section{Evangeliums-Blicke}\label{evangeliums-blicke}

\textbf{BEHARRLICHER GLAUBE.} Diejenigen, die beharrlichen Glauben
zeigen, werden nicht im Unglauben abfallen (Hebr. 3,14). Die
Selbstprüfung in Verbindung mit dem Glauben an das Evangelium bewahrt
vor der Entwicklung eines verhärteten Herzens. Deshalb sollten wir
wachsam sein und an unserem ursprünglichen Vertrauen in das Evangelium
Jesu Christi bis zum Ende festhalten und Gottes Wort dazu benutzen,
jeglichen Unglauben oder Ungehorsam zu offenbaren (Hebr. 3,12-13).
Beharrlichkeit im Glauben ist ein Mittel der Sicherheit, die Gläubige im
christlichen Leben haben (Hebr. 3,6). Glücklicherweise ist es die Gnade
Gottes, die uns rettet und bis zum Ende erhält.

\textbf{SABBATS RUHE.} Im Evangelium gibt Gott seinem Volk die
Sabbatruhe (Hebr. 4,9-10). Im Schöpfungsbericht lesen wir, dass Gott am
siebten Tag ruhte, weil sein Werk vollendet war (1. Mose 2,2). Die
Menschheit ahmt das Muster von Gottes Werk und Ruhe im Sabbatzyklus nach
(Ex 20,8-11; Levitikus 25). Der Sabbat weist nach vorn auf die Ruhe, die
Christus mit seiner Auferstehung und Himmelfahrt erreicht hat (Hebr
10,12-13). Christen treten in eine tiefe Ruhe der Seele ein, wenn sie
bekennen, dass das Erlösungswerk Jesu vollendet war, als er erklärte:
``Es ist vollbracht'', und dann, als er von den Toten auferstand (Joh
19,30; 20,19). Die Kirche wartet nun auf Jesu zweites Kommen, bei dem
sich Gottes Ruhe vollständig offenbaren wird (Offb 22,3-5). Bis zu
diesem Tag sind die Gläubigen aufgerufen, in Christi Erlösung zu ruhen
und der Versuchung zu widerstehen, in die Gerechtigkeit der Werke
zurückzufallen (Gal 2,16; 3,10-14). Ganz-Bibel-Verbindungen

\textbf{TREUER KNECHT.} Mose erfüllte treu die von Gott festgelegten
Aufgaben als Befreier aus der Sklaverei in Ägypten, als Führer des
Exodus und als Gesetzgeber (Ex 3,10; Deut 18,15-19). Jesus ist insofern
größer als Mose, als er sein Volk aus der größeren Sklaverei von Sünde
und Tod befreit und in die verheißene ewige Ruhe Gottes führt.

\textbf{DEN SOHN DES MENSCHEN.} Der Titel, den Jesus mehr als jeder
andere verwendet, um sich auf sich selbst zu beziehen, ist ``der
Menschensohn'' (z.B. Matthäus 8,20; 11,19). Obwohl diese Bezeichnung die
Menschlichkeit Jesu unterstreichen mag, ist der Ausdruck am wichtigsten
in Bezug auf die Gestalt in Daniel 7, die von Gott die höchste Autorität
und ein ewiges Königreich erhält (siehe Dan 7,13-14; Mt 26,64; Mk
14,62). Und als der vollkommene Sohn hält Jesus treu die Gesamtheit des
Gesetzes aufrecht und befähigt ihn, als Hoherpriester und ewiger König
zu dienen (Hebr. 3,1-6).

\textbf{EXODUS ZUR VERHEISSUNG.} Das Volk der Exodus-Generation konnte
aufgrund seines Unglaubens und Ungehorsams nicht in den Rest Gottes (das
verheißene Land) einziehen (Num. 14,20-24; 20,12). Das Muster des Exodus
wiederholt sich in der Befreiung der Kirche von Sünde und Tod und in
ihrer zukünftigen Befreiung von der verheißenen ewigen Ruhe Gottes. Im
Gegensatz zur Exodus-Generation ist Jesus der wahre Israelit, dessen
Treue zu Gott vollkommen ist. Im Gegensatz zu Mose und Josua ist Jesus
in der Lage, jeden einzelnen seines Volkes zu erhalten und der
versprochenen Ruhe Gottes zu übergeben (Hebr 4,1-16).

\section{Theologische Sondierungen}\label{theologische-sondierungen-2}

\textbf{HAUS GOTTES.} Wie andere Metaphern für das Volk Gottes bezieht
sich ``Haus'' auf die gemeinsame Identität der Kirche als Gottes Wohnung
(1 Kor 3,16; 6,19; 2 Kor 6,16; 1 Tim 3,15; 1 Petr 2,5). Jesus hat als
Sohn (Hebr. 3,6), als Baumeister (Hebr. 3,3-4) und als wichtigster
Eckstein (Eph. 2,20-21) einen privilegierten Platz in Gottes Haus. Das
Haus Gottes ist durch den Glauben geeint und aufgefordert, sich
gegenseitig zu Glauben, Gehorsam und Ausdauer zu ermutigen.

\textbf{GOTTES WORT.} Die Bibel ist ein göttlich-menschliches Wort, das
uns von Gott durch Menschen innerhalb der Erlösungsgeschichte gegeben
wurde, die darüber Zeugnis ablegen, wer Gott ist und was er getan hat.
Die Bibel ist Gottes persönliche Äußerung an uns, die als Gott selbst
handelt, unsere Herzen sucht und offenbart. Durch Gottes geschriebenes
Wort kann sein Volk Irrtümer aufdecken, in der Christusähnlichkeit
wachsen und Verständnis dafür gewinnen, was es bedeutet, ein
gottgefälliges Leben zu führen (Hebr 4,12; 2. Tim 3,16).

\section{Persönliche
Implikationen}\label{persuxf6nliche-implikationen-2}

Nehmen Sie sich Zeit, über die Auswirkungen von Hebräer 3,1-4,13 auf Ihr
eigenes Leben heute nachzudenken. Beachten Sie die persönlichen
Implikationen für Ihren Weg mit dem Herrn im Licht der (1)
Evangeliums-Blicke, (2) Gesamtbiblischen Zusammenhänge, (3)
Theologischen Klänge und (4) dieses Abschnitts als Ganzes.

\begin{enumerate}
\def\labelenumi{\arabic{enumi}.}
\tightlist
\item
  Evangeliums-Blicke
\item
  Ganz-Bibel-Verbindungen
\item
  Theologische Sondierungen
\item
  Hebräer 3,1-4,13
\end{enumerate}

Nehmen Sie sich nun einen Moment Zeit, um den Segen und die Hilfe des
Herrn zu erbitten, während Sie mit dem Studium der Hebräer fortfahren.
Und nehmen Sie sich auch einen Moment Zeit, um diese Studieneinheit noch
einmal durchzugehen, über einige Dinge nachzudenken, die der Herr Sie
vielleicht lehrt, und Dinge zu notieren, die Sie in Zukunft noch einmal
überprüfen sollten.

\begin{center}\rule{0.5\linewidth}{0.5pt}\end{center}

Nur für private Zwecke. Übersetzt aus dem Englischen von eurem Diener

Hebrews: A 12-Week Study © 2015 by Matthew Z. Capps. All rights
reserved.

\href{https://www.thegospelcoalition.org/course/knowing-bible-hebrews/\#week-3-jesus-is-superior-to-angelic-beings-heb-15-218}{source}

\bibliography{book.bib,packages.bib}

\end{document}
